\abstracttext{As audiências públicas da Câmara dos Deputados chegam ao público sobretudo por matérias
jornalísticas que selecionam poucas opiniões, sem indicar em que ponto da transcrição cada uma foi
dita. Inspirado nas trilhas de auditoria documental e de mineração de opiniões, este trabalho propõe
um pipeline sobre o PublicHearingBR com três componentes. A Unidade Deliberativa Verificável liga cada
opinião publicada a uma sentença da fala da própria pessoa, com a posição na transcrição, o tipo e a
origem do vínculo; das 2.203 opiniões, 2.105 recebem evidência. Os perfis por ator são escritos por
um modelo de linguagem somente a partir das falas de cada pessoa em audiências de treino, em blocos
de posições, critérios e valores, alinhamentos declarados e forma de argumentar. Na simulação, o
perfil orienta um modelo de linguagem a responder como a pessoa, com trechos recuperados de suas
falas anteriores e uma técnica de decodificação que amplifica a influência desse material
(\textit{classifier-free guidance}); na geração aberta, cada fala é precedida da classificação da
base da estimativa e seguida de uma justificação cujas cópias têm o início conferido
automaticamente. A avaliação usa como gabarito as opiniões publicadas sobre audiências posteriores,
em perguntas de múltipla escolha com distratores da mesma audiência. <Resultados principais: acerto
de cada condição no teste e diferenças entre condições com intervalo de confiança.>}
